\SourceLeaf{601}\Needspace{22\baselineskip}\vspace{6mm}
\pdfbookmark[1]{Resultate aus dem Beobachtungen des magnetischen Vereins im Jahre 1836.}{resultate-1836-anzeige}
\begin{center}\small Gottingische gelehrte Anzeigen.\quad 1837.\quad Juni 29.\end{center}

\textit{Resultate aus den Beobachtungen des magnetischen Vereins im Jahr \textup{1836.} Herausgegeben von \textup{\textsc{Carl Friedrich Gauss}} und \textup{\textsc{Wilhelm Weber}}. Göttingen \textup{1838.} Im Verlag der Dieterichschen Buchhandlung.} (124 Seiten in 8°, dazu 10 Steindrucktafeln.)

Der Titel dieses Werks bezeichnet nur einen Teil seines Inhalts, nämlich denjenigen, der den unmittelbaren Anlass dazu gegeben hat. Von dem Verein, der sich seit mehreren Jahren gebildet hat, um einige der interessantesten Erscheinungen des Erdmagnetismus zu bestimmten vereinbarten Terminen gleichzeitig zu beobachten, war in diesen Blättern schon mehrmals die Rede (1834.\ Aug.\ 9. und 1835.\ März 7.\ [S.\ 518, 529 d.\ B.]). Es ist daher unnötig, Bekanntes hier zu wiederholen. An diesem Verein nimmt bereits eine große Zahl von Orten innerhalb und außerhalb Deutschlands teil, und die Zahl nimmt weiter zu. Die Mitteilung der immer reichhaltiger werdenden Ergebnisse konnte nicht länger auf den privaten Austausch durch Briefwechsel beschränkt bleiben. Für die unmittelbar Beteiligten war die Veröffentlichung im Druck zu einem Bedürfnis geworden. Sie war auch nötig, damit die Ergebnisse zum Gemeingut aller werden können, die sich für die Naturwissenschaften interessieren. Zugleich bietet die Herausgabe die beste Gelegenheit, in besonderen kleineren und größeren Aufsätzen, die mit ihr verbunden werden, nach und nach verschiedene Gegenstände zur Sprache zu bringen. Dazu gehört, was mit dem unmittelbaren Gegenstand direkt zusammenhängt oder indirekt, wie die Instrumente, ihre Berichtigung und Handhabung. aber auch anderes, was in irgendeiner Weise zu einer strengeren wissenschaftlichen Begründung der Lehre vom Magnetismus und Galvanismus beitragen kann.

Die vorliegende erste Lieferung enthält die grafischen Darstellungen der Beobachtungen der magnetischen Variation von sechs Terminen auf ebenso vielen Steindrucktafeln, insgesamt 46 Kurven aus vierzehn verschiedenen Beobachtungsorten, sowie für die letzten drei Termine die Beobachtungen selbst in Zahlen. Den größeren Teil des Werks machen aber außer einer historischen Einleitung
\SourceLeaf{602}die folgenden Aufsätze aus. I. Bemerkungen zur Einrichtung magnetischer Observatorien und Beschreibung der darin aufzustellenden Instrumente. Viele Leser wird es freuen, dass dabei auch die Kosten und Preise angegeben sind. II. Das bei den Beobachtungsterminen anzuwendende Verfahren. III. Auszug aus dreijährigen Beobachtungen der magnetischen Deklination in Göttingen. IV. Beschreibung eines kleinen Apparats für Reisende zur Messung der Intensität des Erdmagnetismus nach absolutem Maß. V. Erläuterungen zu den (hier mitgelieferten) Terminszeichnungen und den Beobachtungszahlen.

Die übrigen vier Steindrucktafeln zeigen einen Lageplan des Göttinger magnetischen Observatoriums, eine perspektivische Darstellung des Beobachtungssaals und der darin aufgestellten Instrumente, seinen Grundriss und genaue Abbildungen aller einzelnen Teile des Magnetometers.

